\documentclass[11pt]{article}
\usepackage[margin=0.9in]{geometry}
\usepackage{amsmath,amssymb,amsthm,booktabs,mathrsfs}
\usepackage[T1]{fontenc}
\usepackage{lmodern,microtype}
\usepackage[colorlinks=true,linkcolor=blue,urlcolor=blue]{hyperref}
\newtheorem{theorem}{Theorem}[section]
\newtheorem{lemma}[theorem]{Lemma}
\newtheorem{proposition}[theorem]{Proposition}
\newtheorem{corollary}[theorem]{Corollary}
\theoremstyle{remark}
\newtheorem{remark}[theorem]{Remark}
\newcommand{\norm}[1]{\left\|#1\right\|}
\title{Toward a four-dimensional Yang--Mills construction\\
\large Notebook 25: a smaller representative of the classical gauge orbit}
\author{Research working document}
\date{17 October 2026}
\begin{document}
\maketitle
\begin{abstract}
The contour-gauge stationary representative need not contract
link angles under blocking. Before applying that rooted current
gauge, however, the partial ordinary average of the established
fine classical minimizer has link-angle bound $C/R$, where
$R=L^s$ is the remaining averaging ratio. We prove the same
factor in its strict rooted support norm and quadratic remainder,
with constants uniform in both averaging depths at fixed base
$L$. The proof applies Banach-valued local analytic estimates
to complete output fibers and does not sum over their increasing
number of coordinates. For an actual one-step gauge-invariant
total correction, this representative supplies an analytic
continuation of its stationary-composition germ wherever the
smaller representative lies in the original domain.
The support-coefficient norm and complete inherited fluctuation
integral still require their own domain estimates.
\end{abstract}

This result addresses the angle-domain obstruction identified in
Notebook 22. It uses the ordinary trajectory bound, the explicit
ordinary minimizer, its localized construction, and the local
gauge identity. At one physical step the displayed amplitude
factor is $C/L$; the constant may depend on that fixed base $L$.
The uniform $C/L^s$ statement for larger $s$ concerns the iterated
constraint, and does not identify a new direct-block integration.

